% Compile with: pdflatex model-diagram.tex
% Then convert to SVG: pdf2svg model-diagram.pdf model-diagram.svg
\documentclass[tikz,border=10pt]{standalone}
\usepackage{tikz}
\usetikzlibrary{shapes,arrows,positioning,fit,calc}
\usepackage{amsmath}
\usepackage{helvet}
\renewcommand{\familydefault}{\sfdefault}
\usepackage{sfmath} % Use sans-serif font for math

\begin{document}

\tikzstyle{input} = [rectangle, rounded corners, minimum width=4cm, minimum height=1.6cm, text centered, draw=black, line width=0.8pt, font=\small, align=center, inner sep=6pt]
\tikzstyle{model} = [rectangle, rounded corners, minimum width=5.8cm, minimum height=4.5cm, text centered, draw=black, line width=1.2pt, font=\normalsize, align=center, inner sep=8pt]
\tikzstyle{output} = [rectangle, rounded corners, minimum width=3.5cm, minimum height=2.7cm, text centered, draw=black, line width=0.8pt, font=\small, align=center, inner sep=6pt]
\tikzstyle{arrow} = [thick,->,>=stealth]
\tikzstyle{outputbox} = [rectangle, rounded corners, fill=white, draw=black, line width=0.5pt, inner sep=2pt]

\begin{tikzpicture}[node distance=1.8cm]

% Title
\node[font=\Large\bfseries, text=black] at (0,9) {Model Architecture: Data Flow};

% ==================== INPUT NODES ====================
\node (input1) [input, fill=white, text=black, text width=3.6cm] at (-8,4.5)
    {\textbf{Tumor Measurements}\\[0.1cm]
    \noindent\rule{3.4cm}{0.4pt}\\[0.1cm]
    \begin{flushleft}\footnotesize
    • SLD \\
    • Time points \\
    • Baseline burden
    \end{flushleft}};

\node (input2) [input, fill=white, text=black, text width=3.6cm] at (-8,1.2)
    {\textbf{Baseline Covariates}\\[0.1cm]
    \noindent\rule{3.4cm}{0.4pt}\\[0.1cm]
    \begin{flushleft}\footnotesize
    • Age, sex, ECOG \\
    • Biomarkers \\
    • Treatment arm
    \end{flushleft}};

\node (input3) [input, fill=white, text=black, text width=3.6cm] at (-8,-2)
    {\textbf{Event Data}\\[0.1cm]
    \noindent\rule{3.4cm}{0.4pt}\\[0.1cm]
    \begin{flushleft}\footnotesize
    • Progression \\
    • Death/censoring
    \end{flushleft}};

% ==================== MECHANISTIC TUMOR DYNAMICS ====================
\node (mech) [model, fill=white, draw=black, line width=1.5pt, text width=5.4cm] at (0,4.5) {
    \large\textbf{Mechanistic Tumor}\\
    \large\textbf{Dynamics}\\[0.1cm]
    \noindent\rule{5.2cm}{0.4pt}\\[0.2cm]
    \small Two-Component State-Space\\[0.3cm]
    \normalsize
    $x_{\text{dec}}(t)$ Decreasing\\[0.12cm]
    $x_{\text{gro}}(t)$ Growing\\[0.3cm]
    \footnotesize\fcolorbox{black}{gray!15}{\parbox{4.8cm}{\centering\textbf{Output:} $\text{SLD}_i(t) \rightarrow T_i^{\text{target}}$}}
};

% ==================== MULTISTATE HAZARD MODEL ====================
\node (events) [model, fill=white, draw=black, line width=1.5pt, text width=5.4cm, minimum height=5.5cm] at (0,-2) {
    \large\textbf{Multistate Hazard Model}\\[0.1cm]
    \noindent\rule{5.2cm}{0.4pt}\\[0.2cm]
    \small Illness-Death Structure\\[0.3cm]
    \begin{flushleft}\footnotesize
    • 0$\to$1: Progression\\
    • 0$\to$2: Death w/o prog.\\
    • 1$\to$2: Post-prog. death\\
    • 0$\to$3: Dropout\\
    • 3$\to$2: Off-trial death
    \end{flushleft}
    \vspace{0.1cm}
    \footnotesize\fcolorbox{black}{gray!15}{\parbox{4.8cm}{\centering\textbf{Output:} $\lambda_{jk}(t) \rightarrow T_i^{01}, T_i^{12}, T_i^{03}$}}
};

% ==================== COMPOSITE ENDPOINTS ====================
\node (pfs) [rectangle, rounded corners, fill=gray!10, draw=black, line width=1.5pt,
             minimum width=4.5cm, minimum height=3.2cm, font=\normalsize, text width=4cm, align=center, inner sep=8pt] at (7.5,1.2) {
    \large\textbf{Composite Endpoints}\\[0.1cm]
    \noindent\rule{3.8cm}{0.4pt}\\[0.15cm]
    \small
    $\text{PFS}_i = \min(T_i^{\text{target}}, T_i^{01})$\\[0.2cm]
    $\text{OS}_i$ via illness-death routing
};

% ==================== OUTPUT NODES ====================
\node (out1) [output, fill=white, text=black, text width=3.2cm] at (13.5,5.5)
    {\textbf{Individual Level}\\[0.1cm]
    \noindent\rule{3cm}{0.4pt}\\[0.1cm]
    \begin{flushleft}\footnotesize
    • Trajectories\\
    • PFS/OS prob.\\
    • RECIST\\
    • PPS
    \end{flushleft}};

\node (out2) [output, fill=white, text=black, text width=3.2cm] at (13.5,1.2)
    {\textbf{Trial Level}\\[0.1cm]
    \noindent\rule{3cm}{0.4pt}\\[0.1cm]
    \begin{flushleft}\footnotesize
    • mPFS, mOS\\
    • PFS/OS rates\\
    • ORR\\
    • KM curves
    \end{flushleft}};

\node (out3) [output, fill=black!90, text=white, text width=3.2cm] at (13.5,-3)
    {\textbf{Decision Support}\\[0.1cm]
    \noindent\textcolor{white}{\rule{3cm}{0.4pt}}\\[0.1cm]
    \small Go/No-Go};

% ==================== ARROWS ====================

% Inputs to Mechanistic
\draw [arrow, black, line width=0.8pt] (input1.east) to[out=0,in=180] (mech.west);
\draw [arrow, black, line width=0.8pt] (input2.east) to[out=0,in=180] (mech.west);

% Inputs to Other Events
\draw [arrow, black, line width=0.8pt] (input2.east) to[out=0,in=180] (events.west);
\draw [arrow, black, line width=0.8pt, dashed] (input3.east) to[out=0,in=180] (events.west);

% Mechanistic to Other Events (SLD feedforward)
\draw [arrow, black, line width=1pt] (mech.south) -- node[right, font=\small, text=black] {$\text{SLD}_i(t)$} (events.north);

% Models to Composite Endpoints
\draw [arrow, black, line width=0.8pt] (mech.east) to[out=0,in=150] node[pos=0.6, outputbox, font=\small] {$T^{\text{target}}$} (pfs.west);
\draw [arrow, black, line width=0.8pt] (events.east) to[out=0,in=210] node[pos=0.6, outputbox, font=\small] {$T^{01}, T^{12}, T^{03}$} (pfs.west);

% PFS to Outputs
\draw [arrow, black, line width=0.8pt] (pfs.north) to[out=90,in=180] (out1.west);
\draw [arrow, black, line width=0.8pt] (pfs.east) -- (out2.west);
\draw [arrow, black, line width=0.8pt] (pfs.south) to[out=270,in=180] (out3.west);

\end{tikzpicture}
\end{document}
